\subsection{李群上的不变点分布场}

定义 4. 设 G 为李群。G 上的一个分布场若在 G 的左（相应地右）平移下不变，则称为左不变的（相应地右不变的）。

换句话说，分布场 \(g \mapsto \Delta_g\)（\(G\) 上的）是左不变的，如果

\[
\Delta_ {g g ^ {\prime}} = \gamma (g) _ {*} \Delta_ {g ^ {\prime}} \quad \text { for } g, g ^ {\prime} \text { in } G,
\]

或者等价地，如果

\[
\Delta_ {g g ^ {\prime}} = \varepsilon_ {g} * \Delta_ {g ^ {\prime}} \quad \text { for   } g, g ^ {\prime} \text {   in   G. }
\]

它是右不变的，如果

\[
\Delta_ {g g ^ {\prime}} = \delta (g ^ {\prime - 1}) _ {*} \Delta_ {g} \quad \text { for   } g, g ^ {\prime} \text {   in   G },
\]

或者等价地，如果.

\[
\Delta_ {g g ^ {\prime}} = \Delta_ {g} * \varepsilon_ {g ^ {\prime}}, \quad \text { for   } g, g ^ {\prime} \text {   in   G. }
\]

定义 5. 设 G 为李群，\(t \in \mathrm{U}(\mathrm{G})\)。用 \(\mathbf{L}_t\) 表示 G 上的分布场 \(g \mapsto \varepsilon_g * t\)，用 \(\mathbb{R}_t\) 表示 G 上的分布场 \(g \mapsto t * \varepsilon_g\)。

换句话说，\(L_{t}\)（相应地 \(R_{t}\)）是由 t 与借助映射 \((g, g') \mapsto gg'\) 在 G 上右（相应地左）作用的 G 所定义的分布场。设 \(\Omega\) 为 G 的开子集，F 为分离多范空间；若 \(f \in \mathcal{C}^{\omega}(\Omega, \mathrm{F})\)，则 \(\mathrm{L}_{t}f = f * t^{\vee} \in \mathcal{C}^{\omega}(\Omega, \mathrm{F})\)，且

\[
\mathrm{R} _ {t} f = t ^ {\vee} * f \in \mathcal {C} ^ {\omega} (\Omega , \mathrm{F})
\]

（第 5 节）。若 G 是有限维的，则微分算子 \(L_{t}\) 与 \(R_{t}\) 是 \(C^{\omega}\) 类的（第 5 节）。

命题 23. (i) 映射 \(t \mapsto \mathbf{L}_t\)（相应地 \(t \mapsto \mathbf{R}_t\)）是向量空间 \(\mathbf{U}(\mathbf{G})\) 到 \(\mathbf{G}\) 上左（相应地右）不变分布所成的向量空间上的同构。

\begin{enumerate}
\def\labelenumi{(\roman{enumi})}
\setcounter{enumi}{1}
\item
  对 U(G) 中的 \(t, t'\)，有 \(L_{t*t'} = L_t \circ L_{t'}\)，\(R_{t*t'} = R_{t'} \circ R_t\)，\(L_t \circ R_{t'} = R_{t'} \circ L_t\)（带有第 5 节的记号滥用）。
\item
  若 \(\theta\) 是映射 \(g \mapsto g^{-1}\)（从 \(\mathbf{G}\) 到 \(\mathbf{G}\) 上），则 \(\theta(\mathbf{L}_t) = \mathbf{R}_{t^{\vee}}\)。
\item
  若 \(t \in \mathrm{U}(\mathrm{G})\) 且 \(g \in \mathrm{G}\)，则 \((\mathrm{L}_t)_g = (\mathrm{R}_{\varepsilon_g * t * \varepsilon_{g^{-1}}})_g\)。
\end{enumerate}

在 G 中，每个右平移与每个左平移交换。于是由第 5 节的命题 21，\(L_{t}\) 是左不变的。由于 \((\mathbf{L}_{t})_{e}=t\)，映射 \(t\mapsto L_{t}\) 是单射的。设 \(\Delta\) 为 G 上的左不变分布场；设 \(t=\Delta_{e}\)；于是 \(\Delta\) 与 \(L_{t}\) 在 e 处有相同的值且都是左不变的，因而 \(\Delta=L_{t}\)。这就对 \(L_{t}\) 证明了 (i)，而对 \(R_{t}\) 的论证是类似的。公式 \(L_{t*t^{\prime}}=L_{t}\circ L_{t^{\prime}}\) 与 \(R_{t*t^{\prime}}=R_{t^{\prime}}\circ R_{t}\) 由 (21) 与 (18) 得出。设 \(t\in\mathrm{U}_{s}(\mathrm{G})\)、\(t^{\prime}\in\mathrm{U}_{s^{\prime}}(\mathrm{G}), f\in\mathcal{C}^{r}(\Omega,\mathrm{F})\)，其中 \(\Omega\) 是 G 中的开集且 \(s+s^{\prime}\leqslant r\)；于是

\[
\begin{array}{r l} \mathrm{L} _ {t} \mathrm{R} _ {t ^ {\prime}} f & = \mathrm{L} _ {t} (t ^ {\prime^ {\vee}} * f) = (t ^ {\prime^ {\vee}} * f) * t \\ & = t ^ {\prime^ {\vee}} * (f * t ^ {\vee}) \quad \text {(Proposition 20)} \\ & = \mathrm{R} _ {t ^ {\prime}} \mathrm{L} _ {t} f \end{array}
\]

因而 \(\mathbf{L}_t\circ \mathbf{R}_{t'} = \mathbf{R}_{t'}\circ \mathbf{L}_t\)。由于 \(\theta\) 是 \(\mathbf{G}\) 到 \(\mathbf{G}^{\vee}\) 的同构，\(\theta (\mathrm{L}_t)\) 是 \(\mathbf{G}\) 上的一个右不变分布场；它在 \(e\) 处的值是 \(\theta^{*}(t) = t^{\vee}\)，因而 \(\theta (\mathrm{L}_t) = \mathrm{R}_{t^{\vee}}\)。最后，

\[
\left(\mathrm{L} _ {t}\right) _ {g} = \varepsilon_ {g} * t = \left(\varepsilon_ {g} * t * \varepsilon_ {g^{-1}}\right) * \varepsilon_ {g} = \left(\mathrm{R} _ {\varepsilon_ {g} * t * \varepsilon_ {g^{-1}}}\right) _ {g}.
\]

注 1. 正是 G 通过右平移在它自身上的作用定义了左不变分布场。

注 2. 设 G 是有限维的。映射

\[
(t, g) \mapsto (\mathbf {R} _ {t}) _ {g} = t * \varepsilon_ {g}
\]

（从 \(\mathrm{U}_{s}(\mathrm{G})\times\mathrm{G}\) 到 \(\mathrm{T}^{(s)}(\mathrm{G})\)）是解析向量丛的同构；因为这个映射是双射的、在每个纤维上是线性的且是解析的（第 5 节）；另一方面，设 \(\phi:\mathrm{T}^{(s)}(\mathrm{G})\to\mathrm{U}_{s}(\mathrm{G})\times\mathrm{G}\) 为其逆双射；若 \(t\in\mathrm{T}_{g}^{(s)}(\mathrm{G})\)，则 \(\phi(t)=(t*\varepsilon_{g^{-1}},g)\)，因而 \(\phi\) 是解析的。同构 \(\phi\) 称为 \(\mathrm{T}^{(s)}(\mathrm{G})\) 的右平凡化。类似地，考虑映射 \((t,g)\mapsto(\mathrm{L}_{\mathrm{t}})_{g}=\varepsilon_{g}*t\)（从 \(\mathrm{U}_{s}(\mathrm{G})\times\mathrm{G}\) 到 \(\mathrm{T}^{(s)}(\mathrm{G})\)）；其逆同构称为 \(\mathrm{T}^{(s)}(\mathrm{G})\) 的左平凡化。通过限制，我们重新得到 \(\mathrm{T}(\mathrm{G})\) 的右平凡化和左平凡化 (§2, 第 2 节)。

